\section{热力学、维里定理与耦合常数积分}

GV §1.8 从 $\varepsilon(r_s)$ 导出压力、压缩率与自旋磁化率，并用维里定理与 Hellmann--Feynman 定理分解动能/势能。这是检验 HF、RPA、QMC 的标准约束。

\subsection{压力、化学势与压缩率}

令 $E=N\varepsilon(r_s)$。在 $d$ 维，
\begin{equation}
\begin{aligned}
  \frac{P}{n}
  &=-\frac{r_s}{d}\varepsilon'(r_s),\\
  \mu
  &=\varepsilon(r_s)-\frac{r_s}{d}\varepsilon'(r_s),\\
  \frac{1}{nK}
  &=\frac{r_s}{d}
  \Bigl[
  \bigl(1-\tfrac{1}{d}\bigr)\varepsilon'(r_s)
  +\frac{r_s}{d}\varepsilon''(r_s)
  \Bigr].
\end{aligned}
\label{eq:PK}
\end{equation}
$K$ 为\textbf{固有}压缩率（proper compressibility）：只对电子气做体积变化、背景视为刚性补偿。实验金属的总压缩率还含离子晶格刚度；固有 $K$ 可在低密变为负值（三维约 $r_s>5.25$），并不表示晶体崩溃。

自旋磁化率 $\chi_S$ 由对极化 $p=(n_\uparrow-n_\downarrow)/n$ 的二阶导数得到。仅含交换时 $K/K_0=\chi_S/\chi_P$；关联使二者分道：压缩率进一步增强，自旋磁化率被压低（GV 图 1.21）。

\subsection{维里定理}

库仑势齐次度为 $-1$，维里定理给出
\begin{equation}
  2t+u
  =d\,\frac{P}{n}
  =-r_s\varepsilon'(r_s),
\label{eq:virial}
\end{equation}
$t$、$u$ 为每粒子动能与势能。结合 $\varepsilon=t+u$，对任意 $d$ 有
\begin{equation}
  t=-\varepsilon-r_s\varepsilon',
  \qquad
  u=2\varepsilon+r_s\varepsilon'.
\end{equation}
更干净的写法是用交换关联能 $\varepsilon_{\mathrm{xc}}=\varepsilon-\varepsilon_0$：
\begin{equation}
  u=2\varepsilon_{\mathrm{xc}}+r_s\varepsilon_{\mathrm{xc}}',
  \qquad
  t=\varepsilon_0-\varepsilon_{\mathrm{xc}}-r_s\varepsilon_{\mathrm{xc}}'.
\end{equation}
相互作用使 $t\ge\varepsilon_0$（动能因关联而升高）。

\subsection{基态能定理（耦合常数积分）}

令电荷标度为 $\sqrt{\lambda}e$（$0\le\lambda\le1$），$H(\lambda)=T+\lambda U$。Hellmann--Feynman 定理
\begin{equation}
  \frac{\mathrm{d}E(\lambda)}{\mathrm{d}\lambda}
  =\bigl\langle\Psi(\lambda)|U|\Psi(\lambda)\bigr\rangle
\end{equation}
给出
\begin{equation}
  E-E_0
  =\int_0^1\frac{\mathrm{d}\lambda}{\lambda}\,
  U(\lambda)
  =N\int_0^1\mathrm{d}\lambda\,
  u(\lambda),
\label{eq:adiab}
\end{equation}
或用结构因子
\begin{equation}
  \varepsilon-\varepsilon_0
  =\frac{1}{2n}
  \int_0^1\mathrm{d}\lambda
  \int\frac{\mathrm{d}^d q}{(2\pi)^d}
  v_{\mathbf{q}}
  \bigl[S_\lambda(\mathbf{q})-1\bigr].
\label{eq:E_from_S}
\end{equation}
这是 RPA 关联能与 DFT 绝热连接公式的共同源头。

\begin{example}[HF 压缩率]
仅取 $\varepsilon=A_d/r_s^2-B_d/r_s$ 代入 \eqref{eq:PK}，三维顺磁气
\begin{equation}
  \frac{K}{K_0}
  =\frac{\chi_S}{\chi_P}
  =\frac{1}{1-r_s/r_s^\star},
  \qquad
  r_s^\star\approx6.03.
\end{equation}
$r_s\to r_s^\star$ 时二者发散，提示 HF 顺磁态的不稳定性（真实转变被关联推迟到低得多的密度）。
\end{example}
